% Generated draft: requires mathematical and source review.
\documentclass[11pt]{article}
\usepackage[margin=25mm]{geometry}
\usepackage{fontspec}
\setmainfont{Latin Modern Roman}
\usepackage{amsmath,amssymb,mathtools,unicode-math}
\setmathfont{Latin Modern Math}
\usepackage{enumitem,needspace,xurl,hyperref}
\setlength{\parindent}{0pt}
\setlength{\parskip}{5pt}
\setlength{\emergencystretch}{3em}
\newcommand{\N}{\mathbb{N}}
\newcommand{\Z}{\mathbb{Z}}
\newcommand{\Q}{\mathbb{Q}}
\newcommand{\R}{\mathbb{R}}
\newcommand{\C}{\mathbb{C}}
\newcommand{\F}{\mathbb{F}}
\newcommand{\A}{\mathbb{A}}
\newcommand{\PP}{\mathbb{P}}
\newcommand{\E}{\mathbb{E}}
\newcommand{\eps}{\varepsilon}
\newcommand{\dd}{\,\mathrm{d}}
\begin{document}
\section*{UnsolvedMath ID: 5500061}
\textbf{Source problem code:} catalogue entry
\subsection{Definitions and mathematical statement}
Let $B_i=\{x\in\R^3:\|x-c_i\|\leq1\}$ for $1\leq i\leq n$; overlaps are explicitly permitted. Count geometric unoriented affine lines once. A line $\ell$ is tangent to $B_i$ when $\operatorname{dist}(c_i,\ell)=1$. Under the closed-ball meaning of ``miss,'' an admissible line has an index set $I\subset\{1,\ldots,n\}$ with $|I|=4$ such that
\[
 \operatorname{dist}(c_i,\ell)=1\quad(i\in I),\qquad
 \operatorname{dist}(c_j,\ell)>1\quad(j\notin I).
\]
Thus it touches exactly four balls and meets no other closed ball. If missing only interiors was intended, additional tangencies and counting multiplicities would need separate treatment.

\textbf{Counting problem.} Determine the number, in the worst case as a function of $n$, of lines tangent to four unit balls and missing all the others. With an intended admissible configuration class $\mathcal C_n$, its extremal form is
\[
 M(n)=\sup_{(B_1,\ldots,B_n)\in\mathcal C_n}
       \#\{\ell:\ell\text{ satisfies the conditions above}\}.
\]
No general-position or isolated-tangent hypothesis is stated. This is essential for a finite combinatorial count: four distinct collinear centers admit a continuous family of parallel common tangent lines, and with $n=4$ there are no other balls to exclude them. The intended restriction, or a convention counting only isolated lines, must be recovered rather than silently added.

The maintained-entry evidence and summary report the problem open and conjecture $\Omega(n^3)$. For a finite-count interpretation that lower-bound statement would mean configurations with at least $cn^3$ admissible lines for all sufficiently large $n$ and a fixed $c>0$; it is not a conjectured upper bound by itself. The summary also mentions up to $12$ common tangents per quadruple, without the nondegeneracy conditions or citation; that finite per-quadruple claim cannot cover arbitrary continuous families. No cubic construction or matching upper bound is asserted here.
\subsection{Short English statement}
Find the worst-case number of lines touching exactly four unit balls in three dimensions and missing all remaining balls, with overlapping balls allowed. The source conjectures a cubic lower bound. Missing general-position or isolated-line conventions are crucial because degenerate configurations can have infinitely many common tangents.
\subsection{Sources}
\begin{itemize}
\item Local UnsolvedMath registry, permanent ID 5500061, source code catalogue entry, ``Lines Tangent to Four Unit Balls''; supplied status open, classification OPEN-TRIAGE, updated 2026-08-06.
\item The Open Problems Project, Computational Geometry, Problem 61, \url{https://topp.openproblem.net/p61}. The record reports source-TeX extraction and access on 2026-07-29; the maintained entry is described as ``Open, conjectured to be $\Omega(n^3)$.''
\item Public source specification: \url{https://github.com/edemaine/topp}. The repository LICENSE is reported to scope software expressly; problem-text reuse remains NEEDS\_REVIEW.
\item The summary mentions a $12$-tangent per-four-sphere result but gives no publication, hypotheses, or theorem reference.
\end{itemize}
\end{document}
